\documentclass[10pt]{article}
\usepackage[T1]{fontenc}
\usepackage[utf8]{inputenc}
\usepackage[french]{babel}
\usepackage{amsmath,amssymb,mathrsfs,mathpazo}
\usepackage{geometry}
\geometry{a4paper,margin=2cm}
\usepackage{graphicx,array,tabularx,tabularray,colortbl,multirow,multicol,booktabs,enumitem,pifont,ulem,fancybox,diagbox,cancel}
\usepackage{pstricks,pstricks-add,pst-plot,pst-tree,pst-node,pst-text,pst-func,pst-math,pst-bspline,pst-3dplot,pst-eucl,pst-solides3d}
\let\pscalculate\undefined
\usepackage{pst-bezier}
\usepackage{tikz,pgfplots,tkz-tab}
\usetikzlibrary{arrows,arrows.meta,patterns,calc,positioning,trees,decorations.markings,shapes.geometric,shapes.misc,shapes.arrows,intersections,decorations.pathmorphing,decorations.pathreplacing,backgrounds,shadows,babel}
\pgfplotsset{compat=1.18}
\usepgfplotslibrary{fillbetween}
\usepackage[np]{numprint}
\usepackage{listings,makecell,pas-tableur,twoopt}
\let\barre\undefined
\usepackage[pstricks]{tablvar}
\definecolor{darkgreen}{RGB}{0,100,0}
\providecommand{\red}{\color{red}}
\providecommand{\blue}{\color{blue}}
\providecommand{\green}{\color{green}}
\providecommand{\black}{\color{black}}
\providecommand{\white}{\color{white}}
\providecommand{\gray}{\color{gray}}
\providecommand{\cyan}{\color{cyan}}
\providecommand{\magenta}{\color{magenta}}
\providecommand{\yellow}{\color{yellow}}
\providecommand{\violet}{\color{violet}}
\providecommand{\orange}{\color{orange}}
\providecommand{\euro}{\texteuro}
\providecommand{\EUR}{\texteuro}
\providecommand{\ieme}{\textsuperscript{e}}
\providecommand{\degres}{\ensuremath{^\circ}}
\providecommand{\degree}{\ensuremath{^\circ}}
\providecommand{\vv}[1]{\overrightarrow{#1}}
\providecommand{\Coord}[1]{#1}
\providecommand{\decosix}{}
\providecommand{\starredbullet}{\ensuremath{\bullet}}
\pagestyle{empty}
\setlength{\parindent}{0pt}
\usepackage{adjustbox}\DeclareUnicodeCharacter{2236}{:}\DeclareUnicodeCharacter{2190}{\ensuremath{\leftarrow}}
\providecommand{\R}{\mathbb{R}}
\providecommand{\N}{\mathbb{N}}
\providecommand{\Z}{\mathbb{Z}}
\providecommand{\Q}{\mathbb{Q}}
\providecommand{\D}{\mathbb{D}}
\providecommand{\C}{\mathbb{C}}
\providecommand{\e}{{\mathrm{e}}}
\providecommand{\vect}[1]{\overrightarrow{#1}}
\providecommand{\vectt}[1]{\overrightarrow{\mathrm{#1}}}
\providecommand{\vvt}[1]{\overrightarrow{\mathrm{#1}}}
\providecommand{\np}[2][]{{#2~#1}}
\providecommand{\vv}[1]{{\overrightarrow{#1}}}
\renewcommand{\d}{{\mathrm{d}}}
\providecommand{\degres}{{^\circ}}
\providecommand{\degre}{{^\circ}}
\providecommand{\blue}{\color{blue}}
\providecommand{\red}{\color{red}}
\providecommand{\enskip}{\quad}
\providecommand{\llbracket}{\lbrack\!\lbrack}
\providecommand{\rrbracket}{\rbrack\!\rbrack}
\providecommand{\euro}{{\text{€}}}
\providecommand{\no}{{\text{n°}}}
\providecommand{\interval}[2]{{[#1;#2]}}
\providecommand{\intervFO}[2]{{[#1;#2[}}
\providecommand{\dotfill}{\ldots}

\providecommand{\ds}{\displaystyle}
\providecommand{\textding}[1]{#1}
\providecommand{\Oij}{\(\mathrm{O};\vec{\imath},\vec{\jmath}\)}
\providecommand{\Oijk}{\(\mathrm{O};\vec{\imath},\vec{\jmath},\vec{k}\)}
\providecommand{\Ouv}{\(\mathrm{O};\vec{u},\vec{v}\)}
\providecommand{\intervFF}[2]{\ensuremath{[#1;#2]}}
\renewcommand{\euro}{\ifmmode\text{\texteuro}\else\texteuro\fi}
\DeclareUnicodeCharacter{00B0}{\ifmmode^{\circ}\else\textdegree\fi}
\DeclareUnicodeCharacter{20AC}{\ifmmode\text{\texteuro}\else\texteuro\fi}
% Énoncés : Éducation nationale. Transcription et corrigés : APMEP, Denis Vergès et contributeurs. Découpage et classement : CoopMaths, Sébastien Lozano, Éric Elter et contributeurs. Adaptation de la présentation : Maths Explorer.
% CC BY-SA : https://coopmaths.fr/www/about/
% https://coopmaths.fr/alea/static/bac/2024/tex/bac_2024_01_sujet0_2.tex
\begin{document}
\noindent\textit{APMEP / CoopMaths. CC BY-SA. Présentation adaptée par Maths Explorer.}\par\medskip
L'exercice est constitué de deux parties indépendantes.


\medskip

\noindent {\Large \textbf{Partie 1}}

\medskip

Pour tout entier $n$ supérieur ou égal à 1, on désigne par $f_{n}$ la fonction définie sur $[0~;~1]$ par :
\[f_{n}(x) = x^{n} \e^{x}\]

On note $C_{n}$ la courbe représentative de la fonction $f_{n}$ dans un repère \Oij{} du plan.

On désigne par $\left(I_{n}\right)$ la suite définie pour tout entier $n$ supérieur ou égal à 1 par :
\[I_{n} = \int_{0}^{1} x^{n} \e^{x} \mathrm{~d} x.\]

\begin{enumerate}
\item 
	\begin{enumerate}
		\item On désigne par $F_{1}$ la fonction définie sur $[0 ; 1]$ par :
\[F_{1}(x) = (x-1) \e^{x}.\]

Vérifier que $F_{1}$ est une primitive de la fonction $f_{1}$.

		\item Calculer $I_{1}$.
	\end{enumerate}
\item À l'aide d'une intégration par parties, établir la relation pour tout $n$ supérieur ou égal à 1 ,
\[I_{n + 1} = \e-(n + 1) I_{n}.\]

\item Calculer $I_{2}$.

\item On considère la fonction \texttt{mystere} écrite dans le langage Python :

\begin{center}
	\renewcommand\arraystretch{1}
	\texttt{\begin{tabular}{l}
			from math import e \# la constante d'Euler e \\
			\\
			def mystere(n):\\
			\quad a = 1\\
			\quad L = [a]\\
			\quad for i in range(1,n):\\
			\quad\quad a = e - (i + 1)*a\\
			\quad\quad L.append(a)\\
			\quad return L\\
	\end{tabular}}
	\renewcommand\arraystretch{1.4}
\end{center}

À l'aide des questions précédentes, expliquer ce que renvoie l'appel \texttt{mystere(5)}.
\end{enumerate}


\medskip

\noindent {\Large \textbf{Partie 2}}

\medskip


\begin{enumerate}
	\item Sur le graphique ci-dessous, on a représenté les courbes $C_{1}, C_{2}, C_{3}, C_{10}, C_{20}$ et $C_{30}$.

\medskip

\hspace*{-1.5cm}	\begin{tikzpicture}
				\begin{axis}[/pgf/number format/use comma,x = {120mm},y = {27mm}, xmin = -0.1, xmax = 1.08, ymin = -0.5, ymax = 2.9, xtick distance = 0.1, ytick distance = 0.5,
					tick label style = {font = \footnotesize},
					 minor tick num = 4, grid = both, major grid style = {line width = .5pt,draw = gray}, axis lines = center]
					\addplot [line width = 1.2pt, color = red, smooth, samples = 100, domain = 0:1 ]{x*exp(x)} node[pos = 0.1,above] {$C_1$};
					\addplot [line width = 1.2pt, color = red, smooth, samples = 100, domain = 0:1 ]{x^2*exp(x)} node[pos = 0.15,above] {$C_2$};
					\addplot [line width = 1.2pt, color = red, smooth, samples = 100, domain = 0:1 ]{x^3*exp(x)} node[pos = 0.2,above] {$C_3$};
					\addplot [line width = 1.2pt, color = red, smooth, samples = 100, domain = 0:1 ]{x^(10)*exp(x)} node[pos = 0.3,left] {$C_{10}~$};
					\addplot [line width = 1.2pt, color = red, smooth, samples = 100, domain = 0:1 ]{x^(20)*exp(x)} node[pos = 0.32,left] {$C_{20}$};
					\addplot [line width = 1.2pt, color = red, smooth, samples = 100, domain = 0:1 ]{x^(30)*exp(x)} node[pos = 0.33,right] {$C_{30}$};
					\node[below left] at (axis cs: 0,0) {{\footnotesize 0}};
				\end{axis}
			\end{tikzpicture}

	\begin{enumerate}
		\item Donner une interprétation graphique de $I_{n}$.

		\item Quelle conjecture peut-on émettre sur la limite de la suite $\left(I_{n}\right)$ ?
	\end{enumerate}
\item Montrer que pour tout $n$ supérieur ou égal à 1,
\[0 \leqslant I_{n} \leqslant \e \int_{0}^{1} x^{n} \mathrm{~d} x.\]
		\item En déduire $\lim\limits_{n \to + \infty} I_{n}$.
	\end{enumerate}

\end{document}
